\documentclass[11pt]{article}
\usepackage[margin=1in]{geometry}
\usepackage[T1]{fontenc}
\usepackage{hyperref}
\title{When Do Simple Baselines Suffice?\\
\large A Frozen-Protocol Simulation Benchmark}
\author{Author list and affiliations require owner verification}
\date{Draft --- not for submission}
\begin{document}
\maketitle

\begin{abstract}
This paper reports a frozen-protocol simulation benchmark of coded tilt,
vibration, acceleration, and multi-node telemetry scenarios. The benchmark
compares simple statistical detectors, unsupervised methods, and an integrated
system using paired seed-level tests across 16 registered scenarios. Its
contributions are a reproducible protocol and paired comparisons, results for
tilt-only conditions, and an identified non-tilt failure mode: the integrated
system gates Isolation Forest positives on tilt or baseline deviation, while
the standalone forest flags some coded non-tilt events. These findings describe
the implemented methods on simulated feature distributions only. No physical
recording, hardware validation, operational trial, or forecasting evaluation
was performed, and the benchmark does not establish a safety benefit.
\end{abstract}

\section{Question and contributions}
The question is when simple baselines suffice within the frozen simulator
protocol, and how the integrated decision rule behaves when an event does not
change tilt. The contributions are:
\begin{enumerate}
  \item A versioned simulation protocol with paired tests across the registered
  16-scenario matrix.
  \item Tilt-only results comparing the current system with simple baselines,
  reported at their fixed code-defined operating points.
  \item A non-tilt failure mode in which the standalone Isolation Forest can
  identify coded feature anomalies but the integrated alert rule suppresses
  them unless a tilt-deviation condition is met.
\end{enumerate}
The benchmark does not claim a new anomaly-detection algorithm or general
superiority over the evaluated baselines.

\section{Related Work}
\textbf{AUTHOR TO COMPLETE: only cite sources read by the author.}

\section{System and simulated data}
The compared methods operate on generated feature sequences from the repository
simulator. The frozen v1 results, raw seed-level metrics, paired comparisons,
and provenance are stored under
\texttt{docs/evaluation/confirmatory\_v1/}. Scenario construction and analysis
decisions are specified in \texttt{docs/evaluation/PROTOCOL.md}; the
interpretation is in \texttt{docs/evaluation/CONFIRMATORY\_ANALYSIS.md}.
No physical sensor recording is included.

\section{Methods}
The frozen protocol compares fixed code-defined operating points and treats a
seed episode as the paired analysis unit. It specifies scenarios, partitions,
event matching, confidence intervals, paired tests, and multiplicity
correction. The methods were not calibrated to a shared false-alarm constraint,
so ranking them does not establish superiority at a common operating point.
The evaluation's \texttt{current\_system} keeps an Isolation Forest positive
only when tilt deviation or baseline deviation reaches its alert threshold;
the standalone \texttt{isolation\_forest} method does not apply that gate.

\section{Results}
The existing confirmatory CSV reports mixed outcomes across the simulated
conditions. On the coded tilt-only cases, simple baselines are competitive
with or outperform the integrated method on event metrics. On selected
non-tilt cases, the standalone forest detects events that the integrated
decision rule often suppresses. The scenario-specific nuisance false-alarm
field is copied from a shared baseline-drift control and is not interpreted as
a per-scenario measurement. No nuisance-alarm rejection conclusion is drawn.
All result claims and their selectors are listed in
\texttt{claims\_ledger.csv}.

\section{Limitations}
All results are from simulation; no physical sensor recording was used. The
scenario distributions are code-defined and are not calibrated to physical
sensor measurements. The integrated system's tilt gate causes missed coded
non-tilt events; changing that behavior requires a new frozen protocol and
future evaluation. The nuisance-rate endpoint cannot support per-scenario
nuisance claims, and thresholds were not calibrated to a shared false-alarm
budget. The benchmark includes no hardware validation, physical deployment,
forecasting evaluation, operational safety assessment, or independent
replication.

\section{Reproducibility and availability}
The protocol hash, environment metadata, raw seed-level metrics, aggregates,
paired tests, unsupported cases, and generated figures are stored under
\texttt{docs/evaluation/confirmatory\_v1/}. Test commands and known
environmental limitations are recorded in \texttt{docs/FINAL\_AUDIT.md}. This
draft is not submission-ready.

\section{Authorship}
Author names, affiliations, contributions, permissions, and acknowledgments
require owner confirmation.

\end{document}
